\section{讨论与局限}

\paragraph{候选排序差异的解释。}
同候选的前 200 条问题上，SC-CL 的 MRR 高于 B1，这与困难候选可能提供更清晰排序信号的设想相容。但这一比较同时改变了候选损失和训练预算，也没有把结构画像与 B1 违例度单独剥离。前缀子集的分布可能不同于剩余测试问题，且全量 B1 闭集分数缺失。因此，目前只能报告配对观察，不能把差异归因于组合采样机制。

\paragraph{排序与生成的取舍。}
SC-CL 在逐候选评分下的 MRR 为 95.27\%，但直接解码 EM 为 87.03\%，略低于 B1 的 87.46\%。两种评估的题目都列出十个选项；差别在于模型自行解码，还是评估程序逐一代入选项评分。因此，两个数值不能互为训练收益证据。如果同预算复现仍呈现排序改善而直接解码不增，方法的用途应限定为候选重排；原始 GRIP 生成任务的整体收益需要单独证实。

\paragraph{过滤的边界。}
train-known 过滤仅排除训练图中该实体对已有的关系。剩余候选可能在持出 split 或真实世界中成立，因此“结构反事实”是训练标签替换的操作名，不是经过验证的假事实。持出正例重叠与开放世界未判定需要分开统计；CoDEx 的实体替换硬负例也不能为本文的关系替换候选提供真值 \citep{safavi2020codex}。若后续审计发现高于随机候选的冲突风险，应相应收窄方法主张或调整训练设计。

\paragraph{成本与适用范围。}
冻结 bank 使候选来源可追溯，但需要离线计算 B1 的全关系分数，也不能随训练中的模型更新。现有资源记录只覆盖 NELL23K 的一次运行，尚不能判断收益是否抵消额外成本。方法针对固定实体对上的关系预测；图级组合问答需要重新定义反事实单位和评估目标。跨图适用性仍取决于 FB15K237-2、WN18RR 和 CoDEx-M 上的同契约评估。
